%% kapa: exam=1 round=27 year=2016 period=2 form=A subject=accounting subject_ko=회계학 seq=40 number=80
%% source: Q-Net 감정평가사 기출문제 (cst00309) artlSeq=5169402 file=extracted/1008/5169402/제27회 감정평가사 1차 문제/제27회 감정평가사 2교시문제(A형).pdf
%% conversion: kordoc 4.24.51 -> tools/gen1.py (auto); review=proofread-sonnet; flags=none
\begin{question}{80}
(주)감평은 A제품을 주문생산하고 있다. 월간 최대 생산가능수량은 10,000개이며, 현재 7,500개를 생산ㆍ판매하고 있다. A제품의 개당 판매가격은 ￦150이며, 현재 조업도 수준하의 원가정보는 다음과 같다.

{\small\centering
\begin{tabular}{|l|r|}\hline
\multicolumn{1}{|c|}{구분} & \multicolumn{1}{c|}{금액} \\ \hline
직접재료원가 & ￦262,500 \\ \hline
직접노무원가 & 300,000 \\ \hline
배치(batch) 수준의 변동원가 & 75,000 \\ \hline
고정제조간접원가 & 275,000 \\ \hline
고정광고비 & 175,000 \\ \hline
\multicolumn{1}{|c|}{계} & 1,087,500 \\ \hline
\end{tabular}\par}

배치 수준의 변동원가는 공정초기화비용(set-up cost), 품질검사비 등을 포함하고 있으며, 1배치에 50개의 A제품을 생산할 수 있다. 최근 (주)감평은 (주)대한으로부터 A제품 2,500개를 개당 ￦120에 구매하겠다는 특별주문을 제안받았다. 이 특별주문을 수락하게 되면 배치를 조정하여 배치당 100개의 A제품을 생산하는 형식으로 변경해야 하고(배치변경에 따른 추가비용은 없음), 기존 고객들에게 개당 ￦10의 할인혜택을 부여해야 한다. 특별주문을 수락한다면 이익에 미치는 영향은?
\begin{choices}
  \item ￦25,000 이익
  \item ￦50,000 이익
  \item ￦25,000 손실
  \item ￦50,000 손실
  \item ￦75,000 손실
\end{choices}
\end{question}
\answer{1}
